% !TeX root = ../Book5.tex
\begin{figure}[htbp]
\centering
\begin{tikzpicture}[x=1cm,y=1cm,every node/.style={font=\small}]
  \draw[AlgebraBlue,line width=.9pt]
    (-2,0) .. controls (-1.3,1.7) and (1.3,1.7) .. (2,0);
  \draw[->,AlgebraBlue] (-.8,1.22)--(-1.1,1.07);
  \draw[->,AlgebraBlue] (.8,1.22)--(1.1,1.07);
  \draw[fill=white] (-2,0) circle[radius=2pt];
  \draw[fill=white] (2,0) circle[radius=2pt];
  \draw[red!70!black] (-2.08,-.08)--(-1.92,.08);
  \draw[red!70!black] (-2.08,.08)--(-1.92,-.08);
  \draw[red!70!black] (1.92,-.08)--(2.08,.08);
  \draw[red!70!black] (1.92,.08)--(2.08,-.08);
  \fill (0,1.275) circle[radius=2pt];
  \node[above] at (0,1.275) {\([1:1:1]\)};
  \node[above] at (-1.8,.35) {\(t\to0\)};
  \node[above] at (1.8,.35) {\(t\to\infty\)};
  \node[below,align=center] at (-2,-.08) {\([0:0:1]\)\\不稳定，删除};
  \node[below,align=center] at (2,-.08) {\([1:0:0]\)\\不稳定，删除};
  \node at (0,.45) {\(y^2=xz\)};
  \begin{scope}[shift={(7,0)}]
    \draw[AlgebraBlue,line width=.9pt] (-2,0)--(2,0);
    \draw[->,AlgebraBlue] (-.8,0)--(-1.2,0);
    \fill (-2,0) circle[radius=2pt];
    \draw[fill=white] (2,0) circle[radius=2pt];
    \draw[red!70!black] (1.92,-.08)--(2.08,.08);
    \draw[red!70!black] (1.92,.08)--(2.08,-.08);
    \fill (0,0) circle[radius=2pt];
    \node[above] at (0,0) {\([1:1:0]\)};
    \node[above] at (-1.2,.25) {\(t\to0\)};
    \node[below,align=center] at (-2,-.08) {\([0:1:0]\)\\半稳定，保留};
    \node[below,align=center] at (2,-.08) {\([1:0:0]\)\\不稳定，删除};
    \node at (0,1.05) {\(z=0\)};
  \end{scope}
\end{tikzpicture}
\caption[射影轨道的边界]{两个射影轨道闭包及其边界点的示意。
左侧轨道的两个边界点都不在半稳定集中，右侧轨道却留下一个半稳定边界点。
空心叉号表示删除的点，实心点表示保留的点；曲线形状仅用于示意闭包中的包含关系。}
\label{b5:fig:projective-torus-boundaries}
\end{figure}
